\chapter{LANDASAN TEORI}
\hyphenpenalty=10000
\exhyphenpenalty=10000
\sloppy

\section{Tinjauan Pustaka}
Penelitian mengenai pengembangan sistem informasi manajemen aset dan peminjaman telah banyak dilakukan untuk mengatasi permasalahan operasional di institusi pendidikan maupun pemerintahan. \textcite{Purwati2022} mengembangkan sistem informasi peminjaman peralatan jaringan dan multimedia menggunakan metode \textit{Waterfall}. Sistem ini dibangun berbasis \textit{website} dengan bahasa pemrograman PHP dan basis data MySQL. Kelebihan penelitian ini adalah keberhasilannya dalam menggantikan pencatatan manual yang rentan hilang, namun penggunaan metode \textit{Waterfall} dinilai kurang fleksibel terhadap perubahan kebutuhan pengguna yang terjadi di tengah proses pengembangan karena sifatnya yang sekuensial.

Dalam konteks manajemen sarana yang lebih luas, \textcite{PutriNabella2025} melakukan penelitian terkait sistem manajemen sarana dan prasarana (Sarpras), IT, serta laboratorium dengan menggunakan metode \textit{Rational Unified Process} (RUP). Sistem ini dikembangkan menggunakan \textit{framework} CodeIgniter 4 dan MySQL untuk memusatkan data aset yang sebelumnya tersebar dalam file Excel. Penelitian ini unggul dalam standarisasi dokumentasi dan integrasi data untuk keperluan audit. Meskipun demikian, pendekatan RUP yang digunakan cenderung berat pada dokumentasi dan kurang adaptif dibandingkan pendekatan \textit{Agile} dalam merespons umpan balik pengguna secara cepat selama fase pengembangan.

Terkait penerapan metodologi yang lebih adaptif, \textcite{Rifky2022} menerapkan kerangka kerja \textit{Scrum} dalam pengembangan sistem informasi pendataan bangunan. Penelitian ini menggunakan PHP dan \textit{framework} CodeIgniter dengan arsitektur MVC. Hasil studi menunjukkan bahwa \textit{Scrum} efektif diterapkan pada proyek dengan batasan waktu yang ketat dan mampu mengakomodasi perubahan fitur yang dinamis melalui mekanisme \textit{sprint}. Sejalan dengan itu, \textcite{Velasquez-Jimenez2025} juga mengimplementasikan perangkat lunak manajemen inventaris menggunakan metodologi \textit{Scrum} yang diuji menggunakan \textit{System Usability Scale} (SUS), yang membuktikan bahwa pendekatan iteratif meningkatkan kepuasan pengguna dan akurasi data inventaris.

\clearpage
Selain aspek metodologi, aspek kinerja teknis \textit{backend} juga menjadi fokus dalam pengembangan sistem modern. \textcite{Pratama2025} melakukan analisis komparatif kinerja RESTful API menggunakan bahasa pemrograman Go, PHP, dan JavaScript dengan metode Raw SQL dan ORM. Temuan penelitian ini menunjukkan bahwa penggunaan bahasa Go dengan metode Raw SQL memberikan kinerja tertinggi dalam menangani beban permintaan dan efisiensi memori dibandingkan PHP. Hal ini mengindikasikan bahwa untuk sistem yang membutuhkan pemrosesan data cepat dan \textit{real-time}, pemilihan teknologi \textit{backend} sangat krusial, aspek yang sering kali kurang diperhatikan dalam penelitian sistem informasi peminjaman sebelumnya yang mayoritas masih berbasis PHP konvensional.

Berdasarkan paparan penelitian terdahulu, mayoritas pengembangan sistem informasi peminjaman sarana prasarana masih didominasi oleh penggunaan metode tradisional seperti \textit{Waterfall} atau RUP, yang memiliki keterbatasan dalam fleksibilitas adaptasi kebutuhan. Di sisi lain, meskipun metode \textit{Scrum} telah diterapkan pada sistem pendataan bangunan dan inventaris, penerapannya secara spesifik pada sistem peminjaman sarana kampus yang mengintegrasikan konsep \textit{CRUD} RDBMS yang optimal, pemrosesan \textit{real-time}, serta prinsip \textit{green computing} masih belum banyak dieksplorasi. Penelitian ini hadir untuk mengisi celah tersebut (\textit{research gap}) dengan mengusulkan pengembangan sistem informasi peminjaman sarana dan prasarana berbasis web menggunakan metode \textit{Scrum} untuk menjamin fleksibilitas pengembangan, serta mengadopsi teknologi yang mendukung efisiensi sumber daya dan kinerja tinggi guna mendukung operasional kampus yang berkelanjutan.

\section{Sistem Informasi dan Manajemen Sarana Prasarana}
Sistem informasi merupakan sistem terintegrasi yang mengombinasikan manusia, perangkat keras, perangkat lunak, jaringan, dan basis data untuk mengolah serta menyajikan informasi guna mendukung operasional dan pengambilan keputusan dalam organisasi \parencite{Herlambang2020}. Dalam konteks perguruan tinggi, sistem informasi menjadi instrumen utama transformasi digital untuk menggantikan proses pengelolaan sarana dan prasarana yang masih manual atau semi-digital, sehingga informasi dapat disajikan secara cepat, akurat, dan konsisten \parencite{Purwati2022}.

Sarana dan prasarana kampus merupakan aset pendukung penting bagi kegiatan akademik dan non-akademik, sehingga pengelolaannya memerlukan manajemen yang sistematis. Proses peminjaman sarana dan prasarana umumnya mencakup tahapan pengajuan, verifikasi ketersediaan, persetujuan, penjadwalan, hingga pengembalian aset. Dalam konteks operasional kampus, proses peminjaman tersebut juga melibatkan peran petugas keamanan sebagai pihak yang memastikan kesesuaian antara jadwal peminjaman pada sistem dengan penggunaan sarana secara aktual di lapangan. Implementasi proses tersebut dalam sistem informasi berbasis web memungkinkan integrasi data inventaris dan transaksi peminjaman secara terpusat serta mendukung pemantauan ketersediaan fasilitas secara \textit{real-time} \parencite{PutriNabella2025}.

Pengelolaan sarana dan prasarana yang masih dilakukan secara manual atau menggunakan \textit{spreadsheet} terpisah berpotensi menimbulkan berbagai permasalahan, seperti redundansi data, keterlambatan pelaporan, dan risiko bentrok jadwal peminjaman akibat keterbatasan informasi ketersediaan fasilitas \parencite{Herlambang2020}. Oleh karena itu, diperlukan sistem informasi peminjaman sarana dan prasarana kampus yang terintegrasi sebagai landasan perancangan sistem pada penelitian ini.

\section{Metodologi \textit{Scrum}}
Pengembangan sistem informasi peminjaman sarana dan prasarana kampus menuntut metode yang mampu beradaptasi terhadap perubahan kebutuhan pengguna selama proses transformasi dari sistem manual ke digital. Berdasarkan penelitian terdahulu, kerangka kerja \textit{Scrum} terbukti efektif dalam pengembangan sistem dengan kebutuhan dinamis dan batasan waktu yang ketat karena mendukung pengembangan bertahap serta validasi berulang pada setiap iterasi \parencite{Rifky2022, Velasquez-Jimenez2025}. Oleh karena itu, \textit{Scrum} dinilai relevan untuk digunakan dalam penelitian ini guna meminimalkan risiko ketidaksesuaian fungsi sistem dengan kebutuhan operasional pengguna.

\textit{Scrum} didefinisikan sebagai kerangka kerja \textit{Agile} yang bersifat iteratif dan inkremental untuk menghasilkan solusi adaptif terhadap permasalahan kompleks melalui transparansi, inspeksi, dan adaptasi \parencite{Schwaber2012}. Dalam penerapannya, \textit{Scrum} melibatkan peran utama \textit{Product Owner}, \textit{Scrum Master}, dan \textit{Development Team}, serta didukung oleh artefak \textit{Product Backlog}, \textit{Sprint Backlog}, dan \textit{Increment} yang dikelola melalui serangkaian event seperti \textit{Sprint Planning}, \textit{Daily Scrum}, \textit{Sprint Review}, dan \textit{Sprint Retrospective} \parencite{permana2015scrum, Nugraha2021}. Karakteristik tersebut menjadikan \textit{Scrum} sebagai landasan metodologis yang sesuai untuk pengembangan sistem peminjaman sarana kampus secara modular dan adaptif, yang pada penelitian ini akan diimplementasikan melalui mekanisme \textit{sprint} pada tahap perancangan dan pengembangan sistem di Bab 3.

\section{Basis Data Relasional dan Konsep \textit{CRUD}}
Sistem peminjaman sarana dan prasarana kampus melibatkan pengelolaan data yang kompleks, mencakup data pengguna, inventaris aset, serta jadwal pemakaian yang bersifat dinamis. Pengelolaan data yang masih dilakukan secara manual atau menggunakan \textit{spreadsheet} terpisah sering menimbulkan permasalahan berupa data yang tidak terintegrasi, kesulitan penelusuran riwayat transaksi, serta keterlambatan penyusunan laporan. Kondisi serupa juga ditemukan pada berbagai institusi pendidikan, di mana keterbatasan sistem pencatatan konvensional berdampak pada rendahnya efisiensi dan akurasi pengelolaan sarana dan prasarana \parencite{PutriNabella2025}.

Pendekatan yang umum digunakan untuk mengatasi permasalahan tersebut adalah penerapan \textit{Relational Database Management System} (RDBMS), yaitu sistem manajemen basis data yang mengorganisasikan data ke dalam tabel-tabel yang saling berelasi melalui kunci utama (\textit{primary key}) dan kunci asing (\textit{foreign key}). Struktur relasional memungkinkan integrasi data secara terpusat, menjaga konsistensi informasi, serta mendukung akses data lintas unit kerja secara \textit{real-time}, sebagaimana diterapkan pada pengembangan sistem manajemen inventaris dan peminjaman di lingkungan pendidikan \parencite{Velasquez-Jimenez2025, Ramdany2024}. Interaksi antara aplikasi dan basis data dalam sistem informasi dilakukan melalui mekanisme \textit{CRUD} (\textit{Create, Read, Update, Delete}), yang merepresentasikan siklus hidup data mulai dari pencatatan aset, pemrosesan peminjaman, hingga pembaruan dan penghapusan data yang tidak lagi relevan \parencite{Purwati2022}.

Untuk menjaga kinerja dan kualitas basis data, diperlukan penerapan teknik normalisasi basis data guna meminimalkan redundansi dan mencegah terjadinya anomali data. Struktur basis data yang ternormalisasi memungkinkan setiap atribut disimpan secara tepat sesuai ketergantungannya, sehingga integritas data tetap terjaga dan efisiensi penyimpanan meningkat. Integrasi antara RDBMS, operasi \textit{CRUD} yang efisien, serta struktur data yang ternormalisasi menjadi landasan teoretis utama dalam perancangan arsitektur data pada penelitian ini, yang implementasi teknisnya akan dibahas lebih lanjut pada Bab 3, sejalan dengan kebutuhan sistem peminjaman yang andal dan berkinerja baik \parencite{Pratama2025}.

\section{Teknologi Sistem Informasi Berbasis Web dengan Prinsip \textit{Green Computing}}
Pengembangan sistem informasi peminjaman sarana dan prasarana kampus memerlukan platform yang mampu mendukung interaksi \textit{multi-user} secara terpusat dengan akses informasi ketersediaan aset secara \textit{real-time} guna meminimalkan bentrok jadwal dan kesalahan pencatatan. Sistem juga harus mengakomodasi akses lintas peran (\textit{role-based access}), mulai dari mahasiswa hingga administrator, untuk menjamin keamanan dan keteraturan alur kerja. Oleh karena itu, teknologi berbasis web dipilih karena mampu mengintegrasikan seluruh proses peminjaman ke dalam satu sistem terpusat yang meningkatkan akurasi data dan efisiensi pelaporan dibandingkan metode manual atau penggunaan \textit{spreadsheet} terpisah \parencite{Herlambang2020, PutriNabella2025, Velasquez-Jimenez2025}.

Secara konseptual, aplikasi web beroperasi menggunakan arsitektur \textit{client–server} yang memungkinkan pemisahan antara antarmuka pengguna dan pemrosesan data di sisi server melalui protokol HTTP. Arsitektur ini menawarkan keunggulan berupa aksesibilitas tinggi, kemudahan integrasi dengan sistem administratif kampus, serta skalabilitas yang baik dalam menangani peningkatan jumlah pengguna \parencite{Purwati2022, Ramdany2024}. Kinerja sistem web sangat bergantung pada komponen \textit{backend} dan pengelolaan basis data dalam menangani logika bisnis serta operasi \textit{CRUD}, di mana efisiensi pemrosesan dan pemilihan metode pengelolaan data yang tepat berperan penting dalam menjaga stabilitas dan waktu respons sistem, terutama pada beban akses yang tinggi \parencite{Pratama2025}.

Sejalan dengan tuntutan keberlanjutan global, pengembangan sistem informasi modern juga harus mengadopsi prinsip efisiensi sumber daya atau \textit{Green Computing}. \textit{Green Computing} atau \textit{Green IT} mencakup praktik perancangan, penggunaan, dan pembuangan sistem komputer dengan dampak lingkungan seminimal mungkin, baik melalui efisiensi energi maupun pengurangan limbah \parencite{Murugesan2012}. Implementasi nyata konsep ini dalam sistem peminjaman adalah penerapan \textit{paperless office}, yang mendigitalkan formulir dan laporan untuk mengurangi konsumsi kertas dan emisi karbon secara signifikan \parencite{Kim2021}. Lebih jauh lagi, prinsip ini juga diterapkan pada level desain perangkat lunak, di mana algoritma dan kode program dioptimalkan untuk meminimalkan beban kerja prosesor dan memori, sehingga mengurangi konsumsi energi operasional server \parencite{Pratama2025}.

\section{Pengujian Perangkat Lunak}
Pengujian perangkat lunak merupakan tahapan fundamental dalam siklus pengembangan sistem untuk menjamin kebenaran, kelengkapan, dan kualitas produk akhir. \textcite{taley2020comprehensive} menjelaskan bahwa pengujian tidak hanya bertujuan untuk menemukan kesalahan (\textit{bugs}), tetapi juga memverifikasi dan memvalidasi bahwa perangkat lunak telah memenuhi persyaratan teknis serta kebutuhan pengguna sebelum dirilis. Dalam konteks pengembangan sistem informasi manajemen, pengujian yang sistematis diperlukan untuk memastikan bahwa alur kerja yang kompleks, seperti proses peminjaman dan validasi data, dapat berjalan tanpa hambatan dan risiko kegagalan fungsi dapat diminimalisir \parencite{Suwirmayanti2020}.

Salah satu metode utama yang digunakan untuk memvalidasi fungsionalitas sistem adalah \textit{Black Box Testing}. Metode ini berfokus pada pengujian spesifikasi fungsional perangkat lunak tanpa perlu mengetahui struktur kode internal atau logika pemrograman di dalamnya \parencite{Suwirmayanti2020}. Pengujian ini dilakukan dengan memberikan serangkaian input pada antarmuka sistem dan mengevaluasi kesesuaian output yang dihasilkan dengan hasil yang diharapkan \parencite{PutriNabella2025}. Pendekatan ini sangat relevan untuk sistem peminjaman sarana kampus guna memastikan bahwa fitur-fitur kritikal, seperti formulir pengajuan, validasi stok aset, autentikasi pengguna, serta skenario validasi penggunaan sarana di lapangan oleh pihak terkait, berfungsi dengan benar sesuai skenario penggunaan (\textit{use case}) yang telah dirancang \parencite{Purwati2022}.

Selain validasi fungsional, evaluasi kualitas teknis pada sisi \textit{backend} dilakukan menggunakan pengukuran \textit{Code Coverage}. \textit{Code Coverage} adalah metrik yang digunakan untuk mengukur seberapa banyak baris kode program yang telah dieksekusi selama proses pengujian otomatis (\textit{unit testing}). Dalam pengembangan \textit{backend} menggunakan bahasa pemrograman Go (Golang), pengujian dan pengukuran coverage didukung secara bawaan oleh pustaka \texttt{testing} dan perintah \texttt{go test} \parencite{kommadi2019learn}. Penggunaan alat ini memungkinkan pengembang untuk mengidentifikasi bagian logika program yang belum teruji (\textit{uncovered code}) dan memastikan bahwa fungsi-fungsi krusial, seperti operasi \textit{CRUD} dan logika bisnis, memiliki tingkat keterujian yang tinggi sebelum integrasi sistem dilakukan.

Kombinasi antara \textit{Black Box Testing} untuk validasi eksternal dan \textit{Code Coverage} untuk validasi internal membentuk strategi pengujian yang komprehensif. Landasan teori ini menjadi dasar bagi pelaksanaan pengujian sistem yang hasil dan analisisnya akan dipaparkan secara mendetail pada Bab 4 Implementasi dan Pengujian, guna membuktikan bahwa sistem informasi peminjaman sarana dan prasarana yang dibangun telah memenuhi standar kualitas yang ditetapkan.
